\documentclass[11pt]{article}
\usepackage[margin=25mm]{geometry}
\usepackage{fontspec}
\setmainfont{Latin Modern Roman}
\usepackage{amsmath,amssymb,mathtools}
\usepackage{unicode-math}
\setmathfont{Latin Modern Math}
\usepackage{xurl,hyperref}
\hypersetup{colorlinks=true,urlcolor=blue}
\setcounter{secnumdepth}{0}
\setlength{\parindent}{0pt}
\setlength{\parskip}{5pt}
\setlength{\emergencystretch}{3em}
\newcommand{\sref}[2]{\hyperlink{source:#1:#2}{[S#2]}}
\newcommand{\cataloguescope}[1]{\paragraph{Recorded target.}#1}
\newcommand{\R}{\mathbb R}
\newcommand{\C}{\mathbb C}
\newcommand{\N}{\mathbb N}
\newcommand{\Z}{\mathbb Z}
\newcommand{\Q}{\mathbb Q}
\newcommand{\F}{\mathbb F}
\newcommand{\eps}{\varepsilon}
\begin{document}
% UnsolvedMath ID 20002535; source catalogue code AIM-PHYSICS-0043
\setcounter{section}{20002534}
\section{Coercivity and collision boundaries for the Broucke-Henon class (1,0,1)}\label{20002535}
{\small UnsolvedMath ID 20002535 · Mathematical Physics}
\subsection{Definitions and mathematical statement}

\paragraph{Shared notation.}$\mathbb N=\{1,2,\ldots\}$; the standard number fields have their usual meanings. All additional conventions are stated locally.


\paragraph{Definitions and notation.}In the Newtonian $n$-body problem, positions $q_i\in\mathbb R^d$ and masses $m_i>0$ obey $m_i\ddot q_i=\nabla_{q_i}U$, where $U(q)=\sum_{i<j}m_im_j/|q_i-q_j|$ and the gravitational constant is normalized to one. Here $T=\tfrac12\sum_i m_i|\dot q_i|^2$, $E=T-U$, and a collision has $q_i=q_j$. The center of mass can be fixed at zero; moment of inertia is $I=\sum_i m_i|q_i|^2$. For a fixed period $P$, the action is $\mathcal A(q)=\int_0^P(T+U)\,dt$. A collision-free minimizer in an admissible loop class satisfies the Newtonian equations when variations are permitted. The period, constraints and closure of the class must be specified: a minimizing sequence can escape or approach collisions. For a planar collision-free periodic loop, identify the plane with $\mathbb C$ and define $k_{ij}=\deg(t\mapsto(q_i(t)-q_j(t))/|q_i(t)-q_j(t)|)\in\mathbb Z$. These pairwise winding numbers describe homology information; they do not determine the full braid homotopy class. The Broucke--Hénon orbit is the named planar periodic family shown in the source. Its class $(1,0,1)$ is the ordered triple of pairwise windings used there. Small and big masses concern mass ratios, not scaling all masses together.

\paragraph{Catalogue wording retained for traceability.}
Prove existence of the Broucke--Hénon orbit for the planar three-body problem. Determine whether such orbits exist for the source’s small-plus-big mass regimes. Determine whether the Broucke--Hénon solution minimizes action in class $(1,0,1)$.
\paragraph{Scope and source qualification.}All three questions remain separate. The source gives numerical and symmetry observations, including failure of coercivity from D2 symmetry alone; these are not an existence or global-minimization proof. The precise unequal-mass parameter range is not supplied in the short question. The catalogue title/status was paired with a machine-generated research attempt; that attempt is not peer-reviewed evidence of settlement. The statement covers the original source questions rather than only the title’s subcase.
\subsection{Short English statement}
Prove Broucke–Hénon periodic orbits, cover unequal small/large mass regimes, and determine minimization in class (1,0,1).
\subsection{Research attempt}
\paragraph{Model, date and investigated scope.}
OpenAI GPT-6 Astra (configured model \texttt{gpt-6-astra}), ultra reasoning; generated 11 October 2026. Proved fixed-positive-mass coercivity for the (1,0,1) winding class with explicit mass dependence. Broucke-Henon existence, singular mass regimes and global minimization remain unproved.

\paragraph{Sources and literature check.}
Read the supplied complete definition, catalogue wording, scope qualifications and source list for this canonical ID. The original-source citations are inherited as supplied references; no fresh primary-source or novelty search was performed for this attempt. No current open-status claim is inferred from this calculation.

\paragraph{Coercivity in the specified winding class.}Take the pair ordering \((k_{12},k_{13},k_{23})=(1,0,1)\), explicitly matching the named pairs for this calculation. Set \(K=\int_0^P\tfrac12\sum_im_i|\dot q_i|^2dt\). For either nonzero-winding pair, weighted Cauchy--Schwarz gives
\[\|\dot q_i-\dot q_j\|_2^2\le2(1/m_i+1/m_j)K.
\]
A nonzero-degree loop \(r=q_i-q_j\) must have its mean within its oscillation radius; otherwise it lies in a disk avoiding zero and has degree zero. Since its oscillation about the mean is at most \(\sqrt P\|\dot r\|_2\),
\[\sup_t|q_i-q_j|\le2\sqrt{2P(1/m_i+1/m_j)K}.
\]
Apply this to pairs 12 and 23, then use \(q_1-q_3=(q_1-q_2)+(q_2-q_3)\). The center-of-mass relation bounds all positions. Thus an action-bounded sequence cannot escape spatially for fixed positive masses.
\paragraph{The mass-ratio and collision limitations.}The constants diverge as a mass tends to zero. They therefore give no uniform compactness across the source's small-plus-big regimes. Moreover finite action permits \(r\sim|t-t_0|^{2/3}\) binary collision behavior; both \(|\dot r|^2\) and \(1/|r|\) are locally integrable. A weak limit can lose its winding class at such a collision even though spatial escape is excluded.
\paragraph{Gap.}A direct minimization proof still needs a collision-exclusion argument and identification of the resulting orbit with the Broucke--Henon geometry. Neither follows from coercivity, and even existence would not settle whether the named orbit globally minimizes the class for every relevant mass ratio. The source's D2-symmetry warning is respected: the estimate here uses the actual winding constraints, not symmetry alone. All three source questions remain unresolved beyond this elementary bound.

\paragraph{Verification and adversarial audit.}
Checked pair ordering, weighted velocity bound, connected winding graph, zero center of mass and blow-up of constants as a mass vanishes. This was an LLM adversarial audit, not independent human verification; no theorem prover was used. Prompts A then B applied to this ID: Checked pair ordering, weighted velocity bound, connected winding graph, zero center of mass and blow-up of constants as a mass vanishes. Scope remains limited: Exclude collisional minimizing limits and identify the resulting orbit, with separate uniform analysis for small-mass regimes.

Coordinator checked the weighted relative-velocity bound, connected pair estimates and failure of uniform control in a vanishing-mass limit. This additional check was model-assisted, not independent human verification.

\paragraph{Tools, process and contributions.}
Python 3 and local shell used to read supplied files and serialize this manuscript; mathematics derived symbolically in the text, with no numerical experiment or formal proof assistant. Contribution prompt A was followed by prompt B. GPT-6 Astra with ultra reasoning drafted the research and performed its stated self-audit; any additional coordinator contribution is identified in the verification and audit. A GPT-6-based Codex coordinator assembled the manuscript and checked repository integration; its exact variant and reasoning configuration were not exposed. No human mathematical contribution, independent human proof check or formal proof certification is supplied. Complete derivations and the consulted source references are included here; private reasoning transcripts are not published.

\paragraph{Final status and first remaining gap.}
\textbf{UNRESOLVED}. Exclude collisional minimizing limits and identify the resulting orbit, with separate uniform analysis for small-mass regimes.

\section{Completion Estimate}
\noindent\textbf{COMPLETION: 0\%}
The calculation is an elementary benchmark or conditional reduction; it establishes no claimed new portion of the complete source target.

\subsection{Sources}
\begin{itemize}
\item[\hypertarget{source:20002535:1}{S1}] ULAM.AI and the UnsolvedMath Contributors, UnsolvedMath, record 20002535 (AIM-PHYSICS-0043). Catalogue curation under CC BY 4.0; definitions and English exposition added locally.
\url{https://huggingface.co/datasets/ulamai/UnsolvedMath}.
\item[\hypertarget{source:20002535:2}{S2}] American Institute of Mathematics, Variational Methods in Celestial Mechanics; notes by Kuo-Chang Chen, version 22 June 2003, source item 6.
\url{https://aimath.org/WWN/varcelest/varcelest.pdf}.
\end{itemize}
\label{end:20002535}
\end{document}
